\documentclass[9pt,letterpaper,onecolumn]{rho-class/rho}
\usepackage{cmap}
\usepackage[spanish,es-nodecimaldot,es-noindentfirst]{babel}
\usepackage{graphicx}
\usepackage{float}
\setlength{\parindent}{0pt}  % Elimina la sangría
\usepackage[utf8]{inputenc}
\usepackage[T1]{fontenc}
\setbool{rho-abstract}{true} % Set false to hide the abstract
\setbool{corres-info}{true} % Set false to hide the corresponding author section
\addbibresource{rho.bib}

%----------------------------------------------------------
% TITLE
%----------------------------------------------------------
\title{Análisis Teórico y Experimental de Algoritmos de Divide y Vencerás: Ordenamiento, Búsqueda y Selección}

%----------------------------------------------------------
% AUTHORS AND AFFILIATIONS
%----------------------------------------------------------
\author[$\dagger$]{Andres Barbosa}
\author[$\dagger$]{Emmanuel Velásquez}
\author[$\dagger$]{Diego Oyarzo}

\affil[$\dagger$]{Universidad de Magallanes}

%----------------------------------------------------------
% DATES
%----------------------------------------------------------
\dates{Este informe fue compilado el \today}

%----------------------------------------------------------
% FOOTER INFORMATION
%----------------------------------------------------------
\smalltitle{Diseño de Algoritmos - Proyecto 2}
\theday{\today} 

%----------------------------------------------------------
% COLORS
%----------------------------------------------------------
\usepackage{xcolor}
\definecolor{codegreen}{rgb}{0,0.5,0}
\definecolor{codegray}{rgb}{0.5,0.5,0.5}
\definecolor{codepurple}{rgb}{0.5,0,0.5}
\definecolor{backcolour}{rgb}{0.95,0.95,0.95}

%----------------------------------------------------------
% Textttt Style
%----------------------------------------------------------
\newtcbox{\mdcode}{on line, tcbox raise base,
  colback=gray!15, colframe=gray!60, boxrule=0.3pt, arc=2pt,
  left=0.3em, right=0.3em, top=0.2ex, bottom=0.2ex, boxsep=0pt,
  before upper=\ttfamily
}
\let\oldtexttt\texttt                                  
\renewcommand{\texttt}[1]{\mdcode{#1}}

%----------------------------------------------------------
% CODE BLOCKS
%----------------------------------------------------------
\lstdefinestyle{CodeStyle}{
    backgroundcolor=\color{backcolour},
    commentstyle=\color{codegreen},
    keywordstyle=\color{codepurple},
    numberstyle=\tiny\color{codegray},
    stringstyle=\color{codepurple},
    basicstyle=\ttfamily\footnotesize,
    breakatwhitespace=false,
    breaklines=true,
    captionpos=b,
    keepspaces=true,
    numbers=left,
    numbersep=5pt,
    showspaces=false,
    showstringspaces=false,
    showtabs=false,
    tabsize=2
}
\lstset{
  literate=
    {á}{{\'a}}1 {é}{{\'e}}1 {í}{{\'i}}1 {ó}{{\'o}}1 {ú}{{\'u}}1
    {Á}{{\'A}}1 {É}{{\'E}}1 {Í}{{\'I}}1 {Ó}{{\'O}}1 {Ú}{{\'U}}1
    {ñ}{{\~n}}1 {Ñ}{{\~N}}1
    {ü}{{\"u}}1 {Ü}{{\"U}}1
}

%----------------------------------------------------------
% PSEUDOCODE
%----------------------------------------------------------
\usepackage{algorithm}
\usepackage[noEnd=true,indLines=true]{algpseudocodex}
\newcommand{\Or}{\textbf{or~}}
\newcommand{\Xor}{\textbf{xor~}}
\newcommand{\To}{\textbf{to~}}
\newcommand{\DownTo}{\textbf{downto~}}
\newcommand{\True}{\textbf{true~}}
\newcommand{\False}{\textbf{false~}}
\newcommand{\Input}{\item[\textbf{Input:}]}
\renewcommand{\Output}{\item[\textbf{Output:}]}
\newcommand{\Print}{\State \textbf{print~}}
\renewcommand{\Return}{\State \textbf{return~}}
\algnewcommand\algorithmicdownto{\textbf{downto~}}
\algnewcommand\Downto{\algorithmicdownto}
\usepackage{amsmath}

%----------------------------------------------------------
% ABSTRACT
%----------------------------------------------------------
\begin{abstract}
El presente proyecto analiza empírica y teóricamente el comportamiento de algoritmos basados en el paradigma de Divide y Vencerás, aplicándolos sobre un robusto sistema de gestión de deportistas construido nativamente en C. El estudio abarca algoritmos de ordenamiento como Merge Sort (en su vertiente clásica y optimizada mediante Insertion Sort) y Quick Sort utilizando el esquema particional de Lomuto con 4 heurísticas de pivote. Asimismo, se incorporan esquemas de Búsqueda Binaria (recursiva y delimitación de rangos), Búsqueda Exponencial y Búsqueda por Interpolación; cerrando el abanico con Quick Select para la búsqueda del k-ésimo elemento. 

Mediante evaluaciones sobre variados tamaños de entrada y arreglos dispares, los resultados arrojan evidencia rigurosa sobre el salto cualitativo en la complejidad temporal asintótica (de $\mathcal{O}(n^2)$ a $\mathcal{O}(n \log n)$) en comparación a técnicas sistemáticas previas y explora en profundidad cómo la selección paramétrica del pivote o los umbrales de combinación moldean determinantemente el rendimiento final en escenarios prácticos.
\end{abstract}

\keywords{Divide y Vencerás, Merge Sort, Quick Sort, Búsqueda Exponencial, Quick Select, Complejidad Temporal, Lenguaje C, Lomuto}

%----------------------------------------------------------
\begin{document}
\maketitle
\thispagestyle{firststyle}
\tableofcontents

\newpage
\section{Introducción}\label{sec:introduccion}

El diseño de algoritmos tiene como objetivo encontrar soluciones a problemas de manera eficiente teniendo en cuenta las limitaciones de tiempo y memoria que existen dentro de los sistemas informáticos.

Tres de los problemas más comunes en la computación son el \textbf{ordenamiento} de datos, la \textbf{búsqueda} de información y la \textbf{selección} de elementos dentro de estructuras de datos; estos problemas se pueden afrontar desde distintas perspectivas, y en este informe los abordamos desde el paradigma de \textbf{Divide y Vencerás}.

El presente informe aborda estos problemas sobre una base de datos ficticia de deportistas (generada de manera procedimental) con el fin de exponer, analizar y comparar algoritmos de ordenamiento (\textit{Merge sort} clásico y optimizado, \textit{Quick sort} con múltiples estrategias de pivote), algoritmos de búsqueda (\textit{Búsqueda binaria recursiva}, \textit{exponencial} e \textit{interpolación}) y selección (\textit{Quick select}). Todo ello analizando su comportamiento a nivel teórico y contrastándolo mediante experimentación empírica.

\section{Objetivo principal}\label{sec:objetivos}
Analizar y comparar el comportamiento de algoritmos de ordenamiento, búsqueda y selección basados en la estrategia de \textit{Divide y Vencerás}, estableciendo una correspondencia entre los resultados teóricos y experimentales aplicados a un sistema de gestión.

\subsection{Objetivos secundarios}\label{sec:objetivos_secundarios}
\begin{enumerate}
    \item Comprender el funcionamiento e implementar correctamente algoritmos con estrategia recursiva de partición.
    \item Analizar empíricamente la complejidad evaluando el impacto de optimizaciones (ej: umbrales en Merge Sort) y variantes (ej: tipos de pivotes en Quick Sort).
    \item Integrar los algoritmos al sistema desarrollado empleando criterios de justificación basados en la eficiencia.
    \item Documentar y contrastar de manera gráfica la mejora sustancial en tiempos de ejecución frente a los algoritmos con complejidad asintótica de orden cuadrático de iteraciones previas.
\end{enumerate}

\newpage
\section{Marco Teórico}\label{sec:marco}

El presente proyecto se apoya en el paradigma de \textit{Divide y Vencerás}, una estrategia de diseño algorítmico que resuelve un problema grande dividiéndolo en subproblemas más pequeños, resolviendo cada uno de forma recursiva y combinando luego sus resultados. Esta técnica suele producir mejoras sustanciales frente a enfoques iterativos o de exploración completa, especialmente cuando las operaciones de división y combinación son más baratas que procesar la entrada completa de manera secuencial.

\subsection{Ordenamiento}
Los algoritmos de ordenamiento estudiados en este trabajo incluyen variantes clásicas y optimizadas de \textit{Merge Sort} y \textit{Quick Sort}. \textit{Merge Sort} garantiza complejidad temporal $\mathcal{O}(n \log n)$ en todos los casos, a costa de memoria auxiliar adicional. \textit{Quick Sort}, por su parte, conserva un comportamiento promedio de $\mathcal{O}(n \log n)$, pero depende de la elección del pivote y puede degradarse a $\mathcal{O}(n^2)$ en escenarios desfavorables.

La versión optimizada de \textit{Merge Sort} introduce un umbral para delegar subarreglos pequeños a \textit{Insertion Sort}, reduciendo el costo constante de la recursión y mejorando el rendimiento empírico sin alterar la cota asintótica principal.

\subsection{Búsqueda}
Para la búsqueda de elementos en arreglos ordenados se emplean técnicas como \textit{Binary Search}, \textit{Exponential Search} e \textit{Interpolation Search}. La búsqueda binaria reduce el espacio de búsqueda a la mitad en cada paso y ofrece complejidad $\mathcal{O}(\log n)$. La búsqueda exponencial permite localizar intervalos rápidamente antes de aplicar una búsqueda binaria, mientras que la interpolación aprovecha la distribución de los datos para estimar la posición probable del elemento.

\subsection{Selección}
\textit{Quick Select} es el algoritmo de selección utilizado para hallar el elemento K-ésimo sin ordenar completamente el arreglo. Su comportamiento promedio es lineal, $\mathcal{O}(n)$, ya que descarta en cada iteración la mitad del espacio que no contiene la respuesta. En el peor caso, al igual que \textit{Quick Sort}, puede alcanzar $\mathcal{O}(n^2)$ si el particionamiento es muy desbalanceado.

\subsection{Criterio de Evaluación}
La comparación experimental del proyecto se basa en medir tiempos de ejecución, observar el efecto de la estructura de los datos y contrastar dichos resultados con el análisis asintótico. De este modo, el marco teórico no solo describe los algoritmos, sino que también justifica por qué ciertas optimizaciones, como el umbral en ordenamiento o la elección del pivote, tienen impacto directo sobre el comportamiento real del sistema.
\newpage
\section{Desarrollo}\label{sec:desarrollo}

\subsection{Análisis Teórico de Algoritmos de Ordenamiento}\label{sec:ordenamiento_dv}
En esta sección se detallan los algoritmos de ordenamiento implementados bajo el paradigma de \textit{Divide y Vencerás}. Se analizará su funcionamiento de manera teórica, desglosando sus pseudocódigos y estudiando su complejidad asintótica global.

\subsubsection{Merge Sort: Clásico y Optimizado}\label{sec:merge_sort}
El algoritmo Merge Sort basa su funcionamiento en dividir el arreglo por la mitad repetidamente hasta obtener arreglos de un solo elemento, para luego fusionarlos (merge) de manera ordenada.
\vspace{0.3cm}
\par \noindent\textit{[Aquí se insertará el pseudocódigo, la explicación teórica y el análisis de la optimización del umbral combinando con Insertion Sort].}

\subsubsection{Quick Sort y sus Variantes de Pivote}\label{sec:quick_sort}
Quick Sort selecciona un elemento como pivote y particiona el arreglo (en este caso implementado según el esquema de partición de Lomuto), ubicando los elementos menores a su izquierda y los mayores a su derecha, para luego procesar los sub-arreglos recursivamente.
\vspace{0.3cm}
\par \noindent\textit{[Aquí se insertará el pseudocódigo, la explicación de las 4 heurísticas de selección de pivote (Último, Primero, Aleatorio, Mediana de tres) y el análisis de su impacto en el mejor, peor y caso promedio].}


\subsection{Análisis Teórico de Algoritmos de Búsqueda}\label{sec:busqueda_dv}
A los algoritmos clásicos de la primera iteración del proyecto se ha sumado un conjunto de algoritmos de búsqueda basados en reducir progresiva o logarítmicamente el espacio restante.

\subsubsection{Búsqueda Binaria: Recursiva y por Rango}\label{sec:busqueda_binaria_rango}
Basada en evaluar permanentemente el valor medio de un sector de estudio asumiendo arreglos ya ordenados, la versión recursiva delega el ciclo a la pila del sistema. Su extensión natural permite no solo atinar a un elemento sino determinar el límite inferior y superior de múltiples coincidencias de este en el arreglo.
\vspace{0.3cm}
\par \noindent\textit{[Aquí se insertará el pseudocódigo, análisis y consideraciones teóricas correspondientes].}

\subsubsection{Búsqueda Exponencial}\label{sec:busqueda_exponencial}
La búsqueda exponencial halla un rango propicio delimitado escalando en potencias de 2, para luego ceder y aplicar una búsqueda binaria local en el tramo hallado.
\vspace{0.3cm}
\par \noindent\textit{[Aquí se insertará el pseudocódigo y la demostración en complejidad analítica].}

\subsubsection{Búsqueda por Interpolación}\label{sec:busqueda_interpolacion}
Mimetizando la forma en que los seres humanos buscan en un diccionario, se realiza una aproximación probabilística del punto objetivo a partir de los extremos del área evaluada para dar atajos rápidos. 
\vspace{0.3cm}
\par \noindent\textit{[Aquí se insertará el pseudocódigo y la complejidad asintótica, profundizando en por qué su peor caso depende directamente de la variabilidad del conjunto].}


\subsection{Análisis Teórico de Algoritmos de Selección}\label{sec:seleccion_dv}

\subsubsection{Quick Select}\label{sec:quick_select}
Basado fuertemente en el esquema de Quick Sort, se detiene prematuramente en la recursividad, desechando la porción del arreglo donde matemáticamente no puede existir el \textit{k-ésimo} elemento buscado.
\vspace{0.3cm}
\par \noindent\textit{[Aquí se insertará el pseudocódigo de manera formal y el análisis de la consecuente complejidad en el mejor y peor de los casos].}


\subsection{Resultados Experimentales}\label{sec:resultados_experimentales}
Esta sección contrasta los hallazgos teóricos de la complejidad con simulaciones obtenidas de los datasets estructurados.

\subsubsection{Impacto del Umbral en Merge Sort}\label{sec:exp_merge_sort_umbral}
\vspace{0.3cm}
\par \noindent\textit{[Aquí se analizarán gráficamente con plots comparativos los distintos umbrales para dirimir cómo rinde frente a la versión clásica pura].}

\subsubsection{Comparativa de Pivotes en Quick Sort}\label{sec:exp_quick_sort_pivotes}
\vspace{0.3cm}
\par \noindent\textit{[Aquí se someterán a prueba, documentando en gráficas, las metodologías de pivotaje estudiadas frente al peor, mejor y caso promedio de los datos de los usuarios].}

\subsubsection{Desempeño en Algoritmos de Búsqueda}\label{sec:exp_busquedas}
\vspace{0.3cm}
\par \noindent\textit{[Aquí irán las observaciones empíricas relacionadas con la agilidad y las caídas bruscas causadas por peores casos medibles].}

\subsubsection{Contraste: Divide y Vencerás frente a Incrementales (Tarea 1)}\label{sec:exp_vs_t1}
\vspace{0.3cm}
\par \noindent\textit{[Aquí se exhibirá en un marco gráfico comparativo la monumental diferencia del desempeño que hay entre O($n^2$) clásicos preexistentes frente a la agilidad del esquema de orden O($n \log n$)].}


\subsection{Aplicación Práctica y Funcionalidades}\label{sec:aplicacion_practica}
\vspace{0.3cm}
\par \noindent\textit{[Aquí se documentará de forma visual y teórica cómo la interfaz ha implementado las nuevas lógicas para listados de rango de puntajes, top atletas, ranking por N criterios o k-ésimo posicionado en el roster central del código fuente. Añadiendo una breve justificación sobre la selección del algoritmo para cada pantalla].}

\newpage
\section{Conclusiones}\label{sec:conclusiones}
A partir del análisis teórico y experimental realizado sobre los algoritmos de la familia \textit{Divide y Vencerás}, y su contraste con los algoritmos iterativos clásicos, se extraen de manera general las siguientes conclusiones:

\subsection*{El salto de cuadrático a log-lineal}
\vspace{0.3cm}
\par \noindent Merge Sort y Quick Sort confirman que el cambio de paradigma hacia Divide y Vencerás reduce de forma marcada el costo de ejecución respecto de los algoritmos cuadráticos estudiados previamente. En las pruebas realizadas, las variantes de ordenamiento log-lineal mantuvieron tiempos estables para tamaños altos de entrada, mientras que los métodos de complejidad $\mathcal{O}(n^2)$ se degradaron de manera visible a medida que creció $N$. Esto valida que la mejora asintótica no es solo teórica, sino que se refleja directamente en la ejecución real.

\subsection*{Optimización empírica: El impacto heurístico}
\vspace{0.3cm}
\par \noindent La experiencia empírica mostró que pequeñas decisiones de implementación cambian de forma tangible el rendimiento final. El umbral de Merge Sort permitió aprovechar Insertion Sort en subarreglos pequeños, disminuyendo el costo de la recursión innecesaria. Del mismo modo, Quick Sort con pivote aleatorio o mediana de tres evitó la degradación observada con pivotes deterministas sobre datos ordenados o invertidos. En consecuencia, la heurística elegida influye directamente en la calidad práctica del algoritmo.

\subsection*{Búsquedas avanzadas}
\vspace{0.3cm}
\par \noindent Las búsquedas binarias, exponenciales e interpolación muestran que un arreglo previamente ordenado abre la puerta a consultas mucho más eficientes que un barrido secuencial. La búsqueda exponencial resulta útil para acotar rápidamente rangos amplios, mientras que la búsqueda por interpolación ofrece buen comportamiento cuando la distribución de los datos es relativamente uniforme. Sin embargo, su efectividad depende fuertemente del perfil de la entrada, lo que confirma que no todas las optimizaciones son universales.

\subsection*{Decisiones arquitectónicas del uso de memoria}
\vspace{0.3cm}
\par \noindent El proyecto también evidenció el costo espacial de cada estrategia. Quick Sort mantiene una ejecución in-place y, por ello, consume menos memoria auxiliar que Merge Sort. En contraste, Merge Sort necesita estructuras temporales para fusionar los subarreglos, a cambio de una mayor estabilidad temporal. La elección entre ambos enfoques depende del equilibrio buscado entre memoria disponible, robustez del tiempo de ejecución y comportamiento frente a distintos tipos de datos.

\subsection*{Cierre de la validación funcional}
\vspace{0.3cm}
\par \noindent La Fase 3 complementó la validación experimental con una revisión estructural del sistema mediante pruebas con 1000 elementos, suficiente para confirmar el flujo completo de generación, lectura, ordenamiento, búsqueda y ejecución de experimentos sin repetir la batería pesada de 25,000 datos usada en la Fase 2. Con ello, el informe queda respaldado tanto por evidencia de rendimiento como por comprobaciones de funcionamiento general de la aplicación.

\newpage
\section{Contribución por integrante}
A continuación se detalla la contribución realizada por cada uno de los integrantes del grupo en el desarrollo del proyecto.
\subsection{Andrés Barbosa}
\textbf{Andrés Barbosa} — Actualizador y mantenedor principal del proyecto.
\begin{itemize}
    \item Responsable de actualizar e integrar el código legado proveniente de su grupo anterior.
    \item Realización de pruebas y validación de las solicitudes del proyecto, asegurando la calidad y funcionalidad del código.
    \item Liderazgo en asignación de tareas y coordinación del equipo, facilitando la comunicación y colaboración entre los integrantes.
\end{itemize}
\subsection{Emmanuel Velásquez}
El estudiante \textbf{Emmanuel Velásquez} contribuyó al proyecto con las siguientes tareas:
\begin{itemize}
    \item Actualización de documentación del repositorio acorde a Proyecto 2 asignado por la docente.
    \item Prototipado del informe, heredando formato anterior y adaptandolo a los requerimientos del proyecto.
    \item Revisión de salud de repositorio, constancia de que mantiene archivos correspondientes al dia para un correcto uso de git.
\end{itemize}

\subsection{Diego Oyarzo}
El estudiante \textbf{Diego Oyarzo} contribuyó al proyecto con las siguientes tareas:
\begin{itemize}
    \item Corrección de errores en análisis de algoritmos.
    \item Inspector de codigo, analisis de algoritmos y pruebas.
\end{itemize}

\newpage
\section{anexos}\label{sec:anexos}
\begin{lstlisting}[language=C, caption={Esquema de partición Lomuto empleado como base en Quick Sort y Quick Select}, label={lst:lomuto}, style=CodeStyle]
/**
 * @brief Esquema de particion de Lomuto
 */
static int lomuto_partition(Player V[], int low, int high, int pivot_type, int (*comp_f)(Player *, Player *)) {
    int p_idx = select_pivot_index(V, low, high, pivot_type, comp_f);
    Player pivot = V[high];
    
    int i = low - 1;
    for (int j = low; j < high; j++) {
        if (comp_f(&V[j], &pivot) <= 0) {
            i++;
            swap_player(&V[i], &V[j]);
        }
    }
    swap_player(&V[i + 1], &V[high]);
    return i + 1;
}
\end{lstlisting}


\newpage
\nocite{*}
\printbibliography
\end{document}
